\begin{lemma}
\label{thm: greedy bias mces policy improves}
Suppose $q$ and $q_+$ satisfy \eqref{eq: definition of q+ greedy section}.
For all policies $\pol \in \grpol(q)$ and $\pol_+ \in \grpol(q_+)$ we have $\pol \preceq \pol_+$ if for all states $s$ where $\pol(s) \neq \pol_+(s)$, the update distribution $\update$ is biased towards the action $\pol(s)$ and
\begin{align}
\label{eq: condition approach below e mcopi proof}
    q(s, \pol_+) \leq q_\pol(s, \pol) .
\end{align}
\end{lemma}
\vspace{-0.8cm}
\begin{proof}
The contraposition of \autoref{thm: cannot exit sliding mode if not all states are better}(1) indicates that
if the update distribution has a greedy bias in state $s$, condition \eqref{eq: condition approach below e mcopi proof} implies that
\begin{align}
    q_{\pol}(s, \pol_+) \geq q_{\pol}(s, \pol) .
\end{align}
As a result, if for all states where the policies differ, condition \eqref{eq: condition approach below e mcopi proof} holds and $\update$ has a greedy bias,
% the advantage is positive in all states where the policies differ and, by 
\autoref{thm: general PIT} ensures that $\pol \preceq \pol_+$.
% policy $\pol_+$ is as good as or better than policy $\pol$.
\end{proof}